%% kapa: exam=1 round=35 year=2024 period=1 form=A subject=economics subject_ko=경제학원론 seq=13 number=53
%% source: Q-Net 감정평가사 기출문제 (cst00309) artlSeq=5213229 file=raw/1008/5213229/1교시 문제지 원본.pdf
%% conversion: kordoc 4.24.51 -> tools/gen1.py (auto); review=proofread-sonnet; flags=none
\begin{question}{53}
비용곡선에 관한 설명으로 옳은 것을 모두 고른 것은?
\begin{framedtext}
ㄱ. 기술진보는 평균비용곡선을 아래쪽으로 이동시킨다. \par
ㄴ. 규모에 대한 수익이 체증하는 경우 장기평균비용곡선은 우하향한다. \par
ㄷ. 단기에서 기업의 평균비용곡선은 한계비용곡선의 최저점에서 교차한다. \par
ㄹ. 규모의 경제가 있으면 평균비용곡선은 수평이다. \par
\end{framedtext}
\begin{choices}
  \item ㄱ, ㄴ
  \item ㄱ, ㄷ
  \item ㄴ, ㄷ
  \item ㄴ, ㄹ
  \item ㄷ, ㄹ
\end{choices}
\end{question}
\answer{1}
